%!TEX root = ../hutbthesis_main.tex

\clearpage
\chapter*{参考文献}
\addcontentsline{toc}{chapter}{参考文献}
\begingroup
\small
\sloppy
\emergencystretch=4em

\begin{enumerate}
\item 伍智锋。分布式飞行仿真技术研究 {[}D{]}. 西安：西北工业大学，2003.
\item 冯聪。无人机集群编队交互式仿真平台的设计与实现 {[}D{]}. 天津：天津大学，2018.
\item 周游。基于 AirSim 仿真平台的无人机三维避障算法研究 {[}D{]}. 成都：电子科技大学，2020.
\item 朱博顺，王成钢，井田。基于 AirSim 平台的无人机蜂群侦察搜索仿真与分析 {[}J{]}. 计算机应用，2021, 41 (S1): 196-201.
\item 郭首江。一种多层无人机集群网络接入控制协议设计及性能优化研究 {[}D{]}. 北京：北京交通大学，2021.
\item 龙腾飞。基于 Airsim 的四旋翼无人机虚实交互平台技术研究 {[}D{]}. 西安：西安工业大学，2023.
\item 冯一飞。基于行为法的分布式无人机集群控制方法与仿真研究 {[}D{]}. 长春：吉林大学，2023.
\item 方璨琦。基于深度强化学习的 AirSim 仿真环境下无人机路径规划研究 {[}D{]}. 广州：华南理工大学，2023.
\item 郑筱宇，胡敏，赵辉宏。基于 AirSim 实现对无人机的位置跟踪及编队控制 {[}J{]}. 齐鲁工业大学学报，2025, 39 (5): 70-75.
\item 章铭泽。无人机集群组网下的通信仿真研究 {[}D{]}. 北京：中国电子科技集团公司电子科学研究院，2025.
\item 方仪豪，邹丹平。多旋翼无人机仿真平台综述 {[}J/OL{]}. 计算机工程，2025: 1-10. https:\allowbreak{}/\allowbreak{}/\allowbreak{}kns.\allowbreak{}cnki.\allowbreak{}net/\allowbreak{}kcms2/\allowbreak{}article/\allowbreak{}abstract?v=3, 2025-01-10.
\item Microsoft Corporation. AirSim: Open Source Simulator for Drones and Autonomous Vehicles{[}EB/OL{]}. https:\allowbreak{}/\allowbreak{}/\allowbreak{}github.\allowbreak{}com/\allowbreak{}microsoft/\allowbreak{}AirSim, 2017-2026.
\item 刘祖均，何明，马子玉。基于分布式一致性的无人机编队控制方法 {[}J{]}. 计 算机工程与应用，2020, 56 (23): 146-152.
\item 李莉，李辉，张雪梅。改进几何控制四旋翼无人机集群轨迹跟踪 {[}J{]}. 兵 器装备工程学报，2025, 46 (3): 232-241.
\item 肖坤，谭帅，王广宇. XTDrone: 可定制化多旋翼无人机仿真平台 {[}C{]}//2020 4th International Conference on Robotics and Automation Sciences. 巴黎: IEEE, 2020: 55-61.
\item 杨浩，孟文，朱书齐。面向无人机集群的数字孪生仿真平台设计 {[}J{]}. 系 统仿真学报，2024, 36 (4): 825-833.
\item 陈勇，林强。多线程调度与线程池在实时仿真系统中的优化方法 {[}J{]}. 计算机工程与设计，2023, 44 (7): 2001-2008.
\item 石晨，张成，刘畅. IMFlySim: 面向无人机集群的高保真仿真平台 {[}C{]}//2021 5th Chinese Conference on Swarm Intelligence and Cooperative Control. 新加坡: Springer, 2023: 209-220.
\item 陈思远。基于 Unity 的无人机集群仿真系统性能瓶颈分析与优化 {[}D{]}. 南京：南京航空航天大学，2022.
\item 吴昊天。实时仿真系统中线程池与 CPU 亲和性调度技术研究 {[}D{]}. 杭州：浙江大学，2023.
\item 林伟，杨玉龙，赵伟。无人机仿真内核态开销过高问题分析与解决 {[}J{]}. 计 算机工程与设计，2022, 43 (8): 2278-2285.
\item 韩冰，曹鹏，孙健。多智能体仿真系统中忙等机制改进方法 {[}J{]}. 控制与决策，2023, 38 (4): 987-994.
\item 高天，徐磊，马晓东。多线程同步机制在无人机实时仿真中的应用比较 {[}J{]}. 电光与控制，2022, 29 (7): 66-71.
\item 董晓阳。面向大规模集群的无人机仿真系统架构设计 {[}D{]}. 哈尔滨：哈尔滨工业大学，2024.
\item 周明宇。基于线程模型优化的无人机仿真系统设计与实现 {[}D{]}. 成都：电子科技大学，2023.
\item 刘畅，陈明，李华。高精度定时器在实时仿真系统中的设计与实现 {[}J{]}. 计算机应用，2022, 42 (S1): 156-161.
\item 朱明亮，邓磊，陈丽。无人机集群仿真中 sleep 与 yield 等待机制性能对比 {[}J{]}. 计算机工程与应用，2024, 60 (3): 167-174.
\item 孙浩，林航，吴昊。基于集中式调度器的多智能体仿真优化方法 {[}J{]}. 系统仿真学报，2023, 35 (9): 1987-1996.
\item 胡俊杰，彭浩，罗强。多线程调度风暴问题分析及工程优化方案 {[}J{]}. 计算机工程，2022, 48 (7): 198-205.
\item 姜浩。基于 AirSim 的多无人机仿真系统性能改进 {[}D{]}. 西安：西北工业大学，2023.
\item 陈涛，杨明，刘军。无人系统仿真中用户态与内核态切换开销研究 {[}J{]}. 计算机应用研究，2023, 40 (4): 1102-1106.
\item 邓宇，陈浩，刘瑞。大规模仿真场景下 GPU 利用率低下问题分析与优化 {[}J{]}. 计算机辅助设计与图形学学报，2022, 34 (11): 1701-1710.
\item 蒋伟，黄涛，李娟。实时仿真系统中内存泄漏与稳定性优化 {[}J{]}. 计算机 工程与设计，2023, 44 (5): 1356-1363.
\end{enumerate}
\endgroup
